% PLANTILLA BEAMER UNIVERSIDAD DE KYUSHU
% Creado por Yu en 2024

\documentclass[11pt]{beamer}
% Puedes usar \documentclass[11pt,aspectratio=169]{beamer} 
% para ajustar la relación de aspecto a 16:9.
\usepackage{float}
\usepackage{booktabs}
\usepackage{graphicx}
\usepackage{xcolor}
\definecolor{vino}{RGB}{128,0,32}
\usepackage{listings}
\usepackage[
	backend=biber,
	style=verbose,
	sorting=ynt
]{biblatex}

\usetheme{Madrid}

\author[Grupo 3]{Catalina Leal Rojas, Lucas Daniel Carrillo Aguirre, Lucas Eduardo Veras Costa}
\title[Sección 1]{Datos}

\setbeamertemplate{navigation symbols}{}

% Descomenta la sección siguiente para mostrar una diapositiva de tabla de contenido
% cada vez que comience una nueva sección
% \AtBeginSection[]
% {
%  \begin{frame}<beamer>
%  \frametitle{Tabla de contenidos}
%  \tableofcontents[currentsection]
%  \end{frame}
% }

\begin{document}

\begin{frame}
\titlepage
\end{frame}

%-----------------------------------------------
\begin{frame}{Fuente de datos}
\small
Los datos utilizados provienen de la \textit{Medición de Pobreza Monetaria y Desigualdad 2018}, elaborada por el \textit{DANE} en el marco de la misión \textit{Empalme de las Series de Empleo, Pobreza y Desigualdad (MESE)}.

\medskip
Esta fuente permite analizar con detalle las condiciones demográficas, laborales, educativas y de vivienda de los hogares colombianos. 

\medskip
Los datos se obtuvieron de la plataforma \textit{Kaggle}, donde se desarrolla una competencia para predecir la pobreza de los hogares.

\medskip
El conjunto se estructura en cuatro archivos: a nivel de hogar e individual, divididos en muestras de \textit{training} y \textit{testing}.
\end{frame}

%-----------------------------------------------
\begin{frame}{Procesamiento y selección de variables}
\small
\begin{itemize}
    \item Se integraron las bases de hogares y personas mediante el identificador \texttt{id}.
    \item Se eliminaron variables con más del 40\% de valores faltantes.
    \item Las variables restantes se seleccionaron según su relevancia teórica y empírica.
    \item Las variables de ingreso del conjunto de entrenamiento no se incluyeron, al no estar disponibles en el conjunto de prueba.
\end{itemize}

\medskip
\textbf{En total:} 26 variables —numéricas, categóricas y binarias— agrupadas en dimensiones laborales, demográficas, educativas y de vivienda.
\end{frame}

%-----------------------------------------------
\begin{frame}{Variables utilizadas}
\small
\begin{table}[H]
\centering
\caption{Descripción y tipo de las variables utilizadas}
\resizebox{\textwidth}{!}{%
\begin{tabular}{lll}
\toprule
\textbf{Variable} & \textbf{Tipo} & \textbf{Descripción} \\
\midrule
Dominio & Categórica & Área metropolitana o zona rural \\
Pobre & Binaria & 1 si el hogar es pobre, 0 de lo contrario \\
Nper & Numérica & Número de personas en el hogar \\
n\_habitaciones & Numérica & Número de habitaciones \\
tipo\_vivienda & Categórica & Tipo de vivienda \\
bin\_headWoman & Binaria & 1 si el jefe del hogar es mujer \\
cat\_educHead & Categórica & Nivel educativo del jefe del hogar \\
t\_ocu & Numérica & Tasa de ocupación \\
t\_inac & Numérica & Tasa de inactividad \\
t\_dependencia & Numérica & Proporción de dependientes en el hogar \\
\bottomrule
\end{tabular}%
}
\end{table}
\end{frame}

%-----------------------------------------------
\begin{frame}{Estadísticas descriptivas (I)}
\small
\begin{itemize}
    \item El conjunto de entrenamiento tiene 164.959 hogares y el de prueba 66.168 ($\approx$40\% de lo de entrenamiento).
    \item Existe un fuerte desbalance: 80\% de hogares no pobres.
\end{itemize}

\medskip
\begin{table}[H]
\centering
\caption{Diferencia de medias entre hogares pobres y no pobres (extracto)}
\small
\begin{tabular}{lcccc}
\toprule
Variable & Pobre & No Pobre & Dif. & Signif. \\
\midrule
N. de Personas & 4.13 & 3.08 & +1.05 & *** \\
N. Habitaciones & 3.03 & 3.48 & -0.45 & *** \\
Edad Jefe & 46.8 & 50.3 & -3.5 & *** \\
Horas Trabajadas Jefe & 28.6 & 34.5 & -5.9 & *** \\
Tasa Ocupación & 0.43 & 0.62 & -0.19 & *** \\
Tasa Inactividad & 0.46 & 0.33 & +0.13 & *** \\
\bottomrule
\end{tabular}
\end{table}
\end{frame}

%-----------------------------------------------
\begin{frame}{Estadísticas descriptivas (II)}
\small
Los hogares pobres presentan:
\begin{itemize}
    \item Más integrantes y mayor número de menores de edad.
    \item Menor nivel educativo y menor proporción con educación superior.
    \item Mayores tasas de inactividad y trabajo por cuenta propia.
\end{itemize}

\medskip
\begin{table}[H]
\centering
\caption{Promedios de variables categóricas}
\small
\begin{tabular}{lcc}
\toprule
\textbf{Variable} & \textbf{Pobre} & \textbf{No Pobre} \\
\midrule
Jefe mujer & 0.47 & 0.41 \\
Jefe ocupado & 0.64 & 0.73 \\
Jefe cuenta propia & 0.47 & 0.32 \\
Vivienda propia total & 0.28 & 0.40 \\
Arriendo & 0.44 & 0.38 \\
Sin título & 0.11 & 0.03 \\
Universitaria (jefe) & 0.10 & 0.32 \\
\bottomrule
\end{tabular}
\end{table}
\end{frame}

%-----------------------------------------------
\begin{frame}{Conclusiones descriptivas}
\small
\begin{itemize}
    \item Los hogares pobres exhiben condiciones laborales más precarias, menor capital humano y mayor dependencia económica.
    \item Las diferencias en educación y tenencia de vivienda son las más marcadas entre grupos.
    \item Estas variables reflejan desigualdades estructurales que el modelo buscará capturar para predecir la pobreza.
\end{itemize}

\medskip
\textbf{Implicación:} las dimensiones laborales, educativas y habitacionales son determinantes esenciales en la clasificación de pobreza.
\end{frame}

%-----------------------------------------------
\end{document}